Configurar o ambiente de desenvolvimento para programar em C é crucial para 
iniciar projetos com eficiência. Abaixo, são fornecidos passos para configurar 
o ambiente em sistemas Linux e Windows.

\subsubsection{Linux}

Para configurar o ambiente de desenvolvimento C em sistemas Linux, siga estes passos:

\begin{enumerate}
    \item \textbf{Instalação do compilador C (GCC):} 
        Em maior parte das distribuições Linux o GCC já está disponível por padrão, 
        mas caso não esteja, ele pode ser facilmente instalado utilizando o gerenciador 
        de pacotes da sua distribuição (como apt para Ubuntu/Debian ou yum para Fedora).
        \\\\
        Segue abaixo uma lista de comandos para instalar o GCC nas distribuições Linux 
        mais famosas:

        \begin{verbatim}
        # Debian
        sudo apt update
        sudo apt install build-essential
        # Red Hat-based (Fedora, CentOS)
        sudo dnf update
        sudo dnf install gcc
        # Arch Linux
        sudo pacman -Syu
        sudo pacman -S gcc
        # openSUSE
        sudo zypper refresh
        sudo zypper install gcc
        # Alpine Linux
        sudo apk update
        sudo apk add build-base

        # Para checar se o GCC foi instalado 
        # corretamente use o comando abaixo
        gcc --version
        \end{verbatim}
\end{enumerate}

\subsubsection{Windows}

Para configurar o ambiente de desenvolvimento C em sistemas Windows, siga 
estes passos:

\begin{enumerate}
    \item \textbf{Instalação do compilador C (GCC):} 
        No Windows a utilização de ferramentas como o GCC acaba por ser 
        um pouco mais complexa pois é necessária a instalação do MinGW ou 
        MSYS2. Essas ferramentas fornecem um ambiente similar ao Linux 
        para compilar código C.

        Para instalar o MinGW, siga os passos abaixo:

        \begin{enumerate}
            \item Acesse o site oficial do MinGW: \texttt{https://osdn.net/projects/mingw/}
            \item Baixe o instalador adequado para o seu sistema (32 ou 64 bits).
            \item Execute o instalador baixado (\texttt{mingw-w64-install.exe}).
            \item Durante a instalação, selecione os pacotes desejados. O 
                pacote essencial para o GCC é o \texttt{mingw32-gcc} para 
                sistemas 32 bits ou \texttt{mingw64-gcc} para sistemas 64 bits.
            \item Siga as instruções do instalador para concluir a instalação.
        \end{enumerate}

        Após a instalação, é necessário configurar o ambiente para usar o GCC via linha de comando:

        \begin{enumerate}
            \item Adicione o diretório bin do MinGW ao caminho do sistema:
                \begin{itemize}
                    \item Para Windows 10, pesquise "Editar variáveis de ambiente 
                        do sistema" na barra de pesquisa e clique em "Variáveis 
                        de ambiente". Em "Variáveis do sistema", selecione a 
                        variável PATH e clique em "Editar...". Adicione o caminho 
                        para o diretório bin do MinGW (exemplo: 
                        \texttt{C:\textbackslash MinGW\textbackslash bin}) e 
                        clique em OK.
                    \item Para versões anteriores do Windows, o processo é 
                        similar, mas pode variar um pouco, caso você utilize 
                        uma versão mais antiga do sistema talvez seja necessário 
                        que você faça uma pesquisa a parte para finalizar o 
                        processo corretamente.
                \end{itemize}
            \item Verifique se a instalação foi bem-sucedida abra o Prompt de Comando (cmd) 
                e insira o seguinte comando:

                \begin{verbatim}
                gcc --version
                \end{verbatim}
        \end{enumerate}

    \item \textbf{Editor de texto ou IDE:} 
        Instale um editor de texto ou uma IDE (Ambiente de Desenvolvimento Integrado), 
        para escrever e compilar código C. Alguns exemplos populares são o Visual 
        Studio Code, Atom, Sublime Text, ou IDEs como o Code::Blocks ou o Eclipse 
        com plugins para C/C++.

        {\tiny Eu pessoalmente recomendo a utilização do Neovim que apesar de ser 
        uma ferramenta de difícil utilização, pode acabar por aumentar muito a sua 
        produtividade no médio/longo prazo.}
\end{enumerate}

Se você seguiu os passos acima com sucesso, seu ambiente de desenvolvimento 
C está configurado e pronto para uso. Agora é o momento de escrever seu 
primeiro programa!
